\chapter{Advanced Features}
\label{ch:advanced-features}

The \textbf{ESP32\_BT\_Controller} firmware is designed as a modular RC + telemetry platform.
Beyond basic joystick control, it includes a lightweight communication protocol, configurable
project modes, a structured telemetry system, and optional I2C expansion.

This chapter highlights the most useful advanced features and explains how they work in practice.

\section{Protocol Overview: Commands, Events, and Telemetry}
\label{sec:protocol-overview}

Communication is packet-based (see \texttt{packets.h}). The firmware uses:

\begin{itemize}
    \item \textbf{RC State packets:} continuous control state (sticks, knobs, switches).
    \item \textbf{Event packets:} short one-shot commands (button presses).
    \item \textbf{Telemetry packets:} feedback from the ESP32 to the Android app (panels, indicators, plots, and configuration labels).
\end{itemize}

All packets use:
\begin{itemize}
    \item Fixed headers (e.g., \texttt{0xCC} followed by a packet type byte)
    \item Packed layouts for predictable parsing
    \item A simple additive checksum (lightweight and fast for embedded systems)
\end{itemize}

\section{Telemetry System (Panels, Indicators, Plots)}
\label{sec:telemetry}

Telemetry is one of the most powerful parts of this firmware. It allows the ESP32 to send
real-time values back to the Android app so the operator can \emph{see what the robot is doing}
while driving it.

\subsection{Telemetry categories}
The firmware generates multiple telemetry streams (implemented in \texttt{telemetry.cpp}):

\begin{itemize}
    \item \textbf{Panels} (left/right numeric values): compact readings such as distance, speed, temperature, or any sensor-derived number.
    \item \textbf{Indicators}: a small set of ``status widgets'' including:
    \begin{itemize}
        \item an analog value (0--100)
        \item a battery level (0--100)
        \item a digital mask (bitfield) for multiple on/off indicators
    \end{itemize}
    \item \textbf{Plots}: fast-updating series for time graphs (default labels include \texttt{Volts}, \texttt{Amps}, and \texttt{RPMs}).
\end{itemize}

\subsection{Telemetry update rates (practical meaning)}
Telemetry is sent periodically (timing constants in \texttt{telemetry.cpp}):

\begin{itemize}
    \item \textbf{Plots:} every 50 ms (\(\approx 20\,Hz\)) --- smooth graphs for fast-changing signals.
    \item \textbf{Indicators:} every 500 ms (\(\approx 2\,Hz\)) --- stable UI values for battery/status.
\end{itemize}

\subsection{Automatic configuration handshake (labels)}
When the phone connects, the ESP32 sends a \textbf{configuration telemetry packet once per connection}
(\texttt{sendConfigTelemetry()}). This packet provides the app with human-readable labels for:

\begin{itemize}
    \item plot series names (e.g., \texttt{Volts}, \texttt{Amps}, \texttt{RPMs})
    \item panel names (e.g., \texttt{Left}, \texttt{Right})
    \item indicator names (e.g., \texttt{Throttle}, \texttt{Battery})
\end{itemize}

This is a major usability feature: the Android app can populate labels automatically without
hardcoding them for each robot build.

\subsection{Panel resend window (robustness feature)}
Panel values use a small resend mechanism (``retry window'') so that when a panel reading changes, it is
retransmitted repeatedly for a short period (see \texttt{resendWindow} and \texttt{resendInterval} in
\texttt{telemetry.cpp}). This improves the chance that the phone receives the updated reading even if a
packet is missed.

\section{Telemetry Sources and Debug Injection}
\label{sec:telemetry-sources}

Telemetry is intentionally split into two layers:

\begin{itemize}
    \item \textbf{\texttt{telemetry.cpp}} is \emph{write-only}: it builds packets and sends them.
    \item \textbf{\texttt{telemetry\_source.cpp}} decides \emph{where data comes from}: hardware reads or debug values.
\end{itemize}

In normal operation, panel/indicator/plot values are read from hardware through \texttt{input\_hw.h}
functions (e.g., numeric inputs and ADC readings). In debug modes, values can be injected over serial
by typing simple comma-separated values. This makes it possible to verify Android UI behavior and packet
decoding \emph{without} real sensors connected.

\paragraph{Why this matters:}
You can develop UI + telemetry displays first, then integrate hardware later, reducing iteration time.

\section{Project Modes: ESP32-Only vs MCP23017 Expansion}
\label{sec:project-modes}

The firmware can be built in different project modes (configured in \texttt{user\_config.h}):

\begin{itemize}
    \item \texttt{MODE\_FULL\_RC}: full RC control using ESP32 GPIO directly.
    \item \texttt{MODE\_FULL\_RC\_MCP}: full RC control using an MCP23017 over I2C for expanded digital I/O.
\end{itemize}

Selecting the correct mode is an ``advanced feature'' because it allows a single codebase to scale from
small robots to larger builds with many switches and indicators.

\section{Using the MCP23017 GPIO Expander (I2C)}
\label{sec:mcp23017}

The MCP23017 adds 16 digital I/O pins over I2C and is supported by the firmware (\texttt{mcp\_io.cpp/.h}).

\subsection{Key details (from the implementation)}
\begin{itemize}
    \item \textbf{I2C address:} \texttt{0x20}
    \item \textbf{Port A (GPIOA):} configured as \textbf{INPUT} (used for reading digital indicators)
    \item \textbf{Port B (GPIOB):} configured as \textbf{OUTPUT} (used for switch and event outputs)
    \item Outputs are tracked in software (\texttt{portBState}) so individual pins can be toggled reliably.
\end{itemize}

\subsection{Practical applications}
\begin{itemize}
    \item A control box with many switches (arm, lights, modes, accessories)
    \item A robot with many status inputs (limit switches, fault lines, bumper sensors)
    \item Complex LED indicator arrays (without consuming ESP32 GPIO)
\end{itemize}

\section{I2C Sensor Support (Examples Included)}
\label{sec:i2c-sensors}

The repository includes example I2C sensor reads in \texttt{i2c\_sensors.cpp}:

\begin{itemize}
    \item A temperature sensor at address \texttt{0x48}
    \item An IMU example at address \texttt{0x68} (accelerometer reads, e.g., MPU6050-style registers)
\end{itemize}

These examples act as templates: you can add initialization steps and additional registers to support
your exact sensor module.

\section{Analog Input Design for Bluetooth Compatibility}
\label{sec:adc1}

A subtle but important advanced design choice is the use of \textbf{ADC1-only} inputs for analog reads
in the default pin configuration. On ESP32, ADC2 is commonly problematic when Wi-Fi/Bluetooth is active,
so the repo pin map keeps analog telemetry on ADC1-capable pins (see \texttt{pins.h}).

\paragraph{Practical implication:}
If you add analog sensors, prefer the configured ADC1 pins (or other ADC1 channels) to avoid intermittent
or invalid readings while Bluetooth is connected.

\section{Feature Flags (Build-Time Modular Firmware)}
\label{sec:feature-flags}

Build-time feature flags are defined in \texttt{feature\_config.h}. They control whether subsystems
are compiled/enabled (PWM outputs, switch outputs, event outputs, telemetry, I2C sensor features, etc.).

\paragraph{Note:}
In the current repository snapshot, \texttt{USE\_TELEMETRY} and \texttt{USE\_I2C\_SENSOR} are defined as
\texttt{0} in both project modes, while telemetry implementation code exists in the repo. This suggests
telemetry can be enabled/disabled by configuration depending on project needs (e.g., to reduce bandwidth or CPU use).

\section{Putting It All Together: Example ``Advanced Build''}
\label{sec:advanced-build-example}

A realistic advanced project combines these features:

\begin{itemize}
    \item \textbf{RC control:} dual sticks + knobs for driving and camera/arm control
    \item \textbf{MCP23017:} many switches and indicators using I2C expansion
    \item \textbf{Telemetry:} battery + current plots, plus panels for temperature and distance
\end{itemize}

This creates a robust operator experience: control remains responsive, while telemetry provides
situational awareness and early warning of battery sag, motor stalls, or sensor faults.